% !TeX program = xelatex
% ============================================================================
%  第 3 课　撤回一步：逆元　—— 教师版讲义
%  与 02-课件/第03课-撤回一步逆元/第03课-撤回一步逆元.tex 同步
% ============================================================================

\jhlesson{3}{撤回一步：逆元}{时长 45 分钟\quad·\quad 教具：无\quad·\quad 本课不发单页}

\begin{mubiao}
  \begin{itemize}
    \item 知道逆元就是能把动作\textbf{撤回}的那一步，会找模 \(12\) 加法的逆元（搭档）。
    \item 会算模 \(7\) 乘法里的逆元，知道这时 \(a\) 叫「可逆的」（也就是能做除法）。
    \item 知道模 \(7\) 乘法里人人有逆元，模 \(6\) 乘法里不是。
  \end{itemize}
\end{mubiao}

\noindent{\small 本课节奏：0—10 回顾（含补上模 12 加法的搭档表）｜10—26 乘 3 的转盘｜26—38 找回头的路｜38—45 收尾。}

\begin{beifen}
  \textbf{时间不够}：若上课 36 分钟还没进模 6 那段，就跳过「一个一个试」那张表，直接报结果
  \(2,4,0,2,4,0\)——没有 \(1\)，因为 \(2x\) 永远是偶数。省下的时间一定留给收尾那句板书：
  「加法里人人都有搭档；乘法里，模 \(7\) 有，模 \(6\) 没有。」
  \textbf{时间富余}：让学生再试一个模 \(5\) 的乘法 \(2x\equiv1\pmod5\)（答案 \(3\)），
  看到「换成质数就又人人有逆元了」，为第 11、12 课埋一层伏笔。
\end{beifen}

\section*{一、回顾（0—10 分钟）}

\begin{caozuo}
  手指回答：模 \(12\) 加法里，\(3\)、\(7\)、\(5\) 的搭档分别是谁？——\(9\)、\(5\)、\(7\)。
\end{caozuo}

再把上节课的两件事摆出来：模 \(12\) 加法的单位元是 \(0\)，每个元素都有搭档；
有理数乘法的单位元是 \(1\)，非零数都有倒数。

\noindent 上节课最后留了一个空（模 \(12\) 加法那一行），现在补上：

\begin{center}
{\small
\begin{tabular}{@{}l*{12}{c}@{}}
  \toprule
  \(a\) & 0 & 1 & 2 & 3 & 4 & 5 & 6 & 7 & 8 & 9 & 10 & 11 \\
  \midrule
  搭档 & 0 & 11 & 10 & 9 & 8 & 7 & 6 & 5 & 4 & 3 & 2 & 1 \\
  \bottomrule
\end{tabular}}
\end{center}

\(1\) 配 \(11\)（\(1+11=12\equiv0\)），\(2\) 配 \(10\)，\(3\) 配 \(9\)，……

\begin{zhu}
  有两个数只配自己——\(0\)（\(0+0=0\)）和 \(6\)（\(6+6=12\equiv0\)）。
  这一条要专门指出来，第 7 课数「阶」的时候还会碰到它。
\end{zhu}

一句话过渡：加法里的「回头」我们已经会了，今天换成\textbf{乘法}。

\section*{二、乘 3 的转盘（10—26 分钟）}

在模 \(7\) 的世界里做乘法，只看 \(1,2,3,4,5,6\)。

\begin{zhu}
  把 \(0\) 请出去的理由要先讲清楚：\(0\) 乘任何数都得 \(0\)，永远回不到 \(1\)，
  它没有逆元。今天研究的正是逆元，所以先把它放在一边。
\end{zhu}

板书两行（一边写一边说）：

\begin{center}
\begin{tabular}{@{}lcccccc@{}}
  \toprule
  原来的数 & \(1\) & \(2\) & \(3\) & \(4\) & \(5\) & \(6\) \\
  \midrule
  乘 \(3\) 之后 & \(3\) & \(6\) & \(2\) & \(5\) & \(1\) & \(4\) \\
  \bottomrule
\end{tabular}
\end{center}

算式逐个报出来：\(3\times1=3\)，\(3\times2=6\)，\(3\times3=9\equiv2\)，
\(3\times4=12\equiv5\)，\(3\times5=15\equiv1\)，\(3\times6=18\equiv4\)。

\begin{caozuo}
  全班一起念一遍：「\(1\) 变成 \(3\)，\(2\) 变成 \(6\)，\(3\) 变成 \(2\)，\(4\) 变成 \(5\)，
  \(5\) 变成 \(1\)，\(6\) 变成 \(4\)。」念完停下来问：
  \textbf{有没有两个数变成同一个？}——没有，\(3,6,2,5,1,4\) 六个数字一个不多一个不少。
\end{caozuo}

所以要重说一遍：\textbf{「乘 \(3\)」这个动作，把 \(1,2,\dots,6\) 六个位置重新排了一遍。}

\begin{zhu}
  这一步要慢，不要赶。它是第 12 课费马小定理证明的伏笔：
  「乘 \(3\)」只是把整张表重排，一个位置也没压扁。
\end{zhu}

\begin{caozuo}
  手指回答：谁乘 \(3\) 之后变成 \(6\)？（\(2\)）谁乘 \(3\) 之后变成 \(1\)？（\(5\)）
  \(5\) 乘 \(3\) 之后变成几？（\(1\)）
\end{caozuo}

\section*{三、找回头的路（26—38 分钟）}

抛问题：\textbf{谁是 \(3\) 的逆元？}也就是，什么数 \(x\) 满足 \(3x\equiv1\pmod7\)？

在加法里「回头」是加上相反数；在乘法里，「回头」就是乘一个数，让结果回到单位元 \(1\)。

\begin{caozuo}
  当着全班板演，把 \(x\) 从 \(1\) 试到 \(6\)，把结果填进表里：
  \[
  \begin{array}{@{}lcccccc@{}}
    \toprule
    x & 1 & 2 & 3 & 4 & 5 & 6 \\
    \midrule
    3x \bmod 7 & 3 & 6 & 2 & 5 & \mathbf{1} & 4 \\
    \bottomrule
  \end{array}
  \]
  走到 \(x=5\) 时停下：\(3\times5=15=2\times7+1\equiv1\pmod7\)。
  所以 \textbf{\(3\) 的逆元是 \(5\)}：乘 \(3\) 的逆动作，就是乘 \(5\)。
\end{caozuo}

\begin{definition}[可逆]
  设 \(e\) 是运算 \(*\) 的单位元。如果元素 \(a\) 有逆元，
  就说 \(a\) 在这种运算下是\textbf{可逆的}。
\end{definition}

换一个模数再试。问：\(2\) 在模 \(6\) 的乘法下有逆元吗？也就是解 \(2x\equiv1\pmod6\)。

\begin{caozuo}
  同样从 \(1\) 试到 \(6\)：
  \[
  \begin{array}{@{}lcccccc@{}}
    \toprule
    x & 1 & 2 & 3 & 4 & 5 & 6 \\
    \midrule
    2x \bmod 6 & 2 & 4 & 0 & 2 & 4 & 0 \\
    \bottomrule
  \end{array}
  \]
  表里根本没有 \(1\)。原因一看就明白：\(2x\) 永远是偶数，永远不可能是 \(1\)。
  更根本地说，\(2\) 和 \(6\) 有公因数 \(2\)。
\end{caozuo}

\section*{四、收尾（38—45 分钟）}

对照表：

\begin{center}
\begin{tabular}{@{}lcc@{}}
  \toprule
   & \textbf{单位元} & \textbf{逆元} \\
  \midrule
  有理数加法 & \(0\) & 相反数 \(-a\) \\
  有理数乘法 & \(1\) & 倒数 \(1/a\ (a\neq0)\) \\
  模 \(7\) 乘法 & \(1\) & 表里的搭档（\(1\leftrightarrow1,\ 2\leftrightarrow4,\ 3\leftrightarrow5,\ 6\leftrightarrow6\)） \\
  \bottomrule
\end{tabular}
\end{center}

验一验：\(2\times4=8\equiv1\)，\(3\times5=15\equiv1\)，\(6\times6=36\equiv1\)。

板书收尾那句话：\textbf{加法里人人都有搭档；乘法里，模 \(7\) 有，模 \(6\) 没有。}

\begin{zhu}
  有理数乘法里 \(0\) 没有倒数，和模 \(6\) 里的 \(2\) 一样，是被同一件事挡在门外的。
  这句话只说现象，不下结论，留到后面几节课再说清楚。
\end{zhu}

\begin{xiaojie}
  \begin{itemize}
    \item 逆元是把动作\textbf{撤回}的那一步：加上 \(a\) 用 \(-a\) 撤回，乘上 \(3\) 用乘 \(5\) 撤回。
    \item 在模质数的乘法里，非零元素都有逆元，所以能做除法。
    \item 模 \(6\) 的乘法里，\(2\) 没有逆元——它有公因数挡路。
  \end{itemize}
\end{xiaojie}

\noindent\textbf{下节课}：把这几套共同遵守的规矩装进一个定义——群。